\textbf{Phần 1. Trắc nghiệm lựa chọn}
\setcounterpageref{ex}{0}

\begin{ex} %[3D11N59]
Số gồm $3$ chục nghìn, $4$ nghìn, $5$ trăm, $2$ chục và $7$ đơn vị viết là
\choice
{\True $34527$}
{$34257$}
{$35427$}
{$43527$}
\loigiai{
Số đó là $34527$.\\
\textbf{Đáp án của câu hỏi là A}
}
\end{ex}

\begin{ex} %[3D11N59]
Số $100\,000$ đọc là
\choice
{Mười nghìn}
{\True Một trăm nghìn}
{Một triệu}
{Một vạn}
\loigiai{
Số $100\,000$ đọc là Một trăm nghìn.\\
\textbf{Đáp án của câu hỏi là B}
}
\end{ex}

\begin{ex} %[3D11H59]
Trong số $74\,852$, chữ số 7 có giá trị là
\choice
{$7000$}
{$700$}
{$70$}
{\True $70000$}
\loigiai{
Chữ số 7 ở hàng chục nghìn nên có giá trị là $70\,000$.\\
\textbf{Đáp án của câu hỏi là D}
}
\end{ex}

\begin{ex} %[3D11H59]
Số liền sau của số $99\,999$ là
\choice
{\True $100000$}
{$99998$}
{$10000$}
{$99990$}
\loigiai{
$99999 + 1 = 100\,000$.\\
\textbf{Đáp án của câu hỏi là A}
}
\end{ex}

\begin{ex} %[3D11H59]
Số liền trước của số $50\,000$ là
\choice
{$50001$}
{$49990$}
{\True $49999$}
{$49000$}
\loigiai{
$50000 - 1 = 49999$.\\
\textbf{Đáp án của câu hỏi là C}
}
\end{ex}

\begin{ex} %[3D11H59]
Số \textit{Sáu mươi nghìn không trăm linh năm} viết là
\choice
{$60500$}
{\True $60005$}
{$60050$}
{$65000$}
\loigiai{
Số đó viết là $60005$.\\
\textbf{Đáp án của câu hỏi là B}
}
\end{ex}

\begin{ex} %[3D11V59]
Số gồm $8$ chục nghìn, $5$ trăm và $3$ đơn vị viết là
\choice
{\True $80503$}
{$85003$}
{$85030$}
{$80530$}
\loigiai{
Số đó là $80503$.\\
\textbf{Đáp án của câu hỏi là A}
}
\end{ex}

\begin{ex} %[3D11V59]
Viết số $54\,230$ thành tổng các hàng:
\choice
{$5000 + 4000 + 200 + 30$}
{$50000 + 400 + 200 + 30$}
{\True $50000 + 4000 + 200 + 30$}
{$54000 + 230$}
\loigiai{
$54230 = 50000 + 4000 + 200 + 30$.\\
\textbf{Đáp án của câu hỏi là C}
}
\end{ex}

\begin{ex} %[3D11V59]
Từ 5 chữ số $1, 0, 3, 5, 8$ lập được số lớn nhất có 5 chữ số khác nhau là
\choice
{$85130$}
{$85301$}
{$58310$}
{\True $85310$}
\loigiai{
Số lớn nhất là $85310$.\\
\textbf{Đáp án của câu hỏi là D}
}
\end{ex}

\begin{ex} %[3D11V59]
Từ 5 chữ số $2, 0, 4, 6, 9$ lập được số bé nhất có 5 chữ số khác nhau là
\choice
{$24069$}
{\True $20469$}
{$20496$}
{$02469$}
\loigiai{
Số bé nhất là $20469$.\\
\textbf{Đáp án của câu hỏi là B}
}
\end{ex}

\begin{ex} %[3D11V59]
Dãy số tròn chục nghìn gồm các số:
\choice
{$1000, 2000, 3000, \dots$}
{$100, 200, 300, \dots$}
{\True $10000, 20000, 30000, \dots, 90000$}
{$10000, 11000, 12000, \dots$}
\loigiai{
Các số tròn chục nghìn là $10000, 20000, \dots, 90000$.\\
\textbf{Đáp án của câu hỏi là C}
}
\end{ex}

\begin{ex} %[3D11C59]
Tìm số có 5 chữ số biết chữ số hàng chục nghìn là 4, chữ số hàng nghìn gấp đôi chữ số hàng chục nghìn, các chữ số còn lại đều bằng 0.
\choice
{$44000$}
{$46000$}
{$84000$}
{\True $48000$}
\loigiai{
Chục nghìn: 4. Hàng nghìn: $4 \times 2 = 8$. Số đó là $48000$.\\
\textbf{Đáp án của câu hỏi là D}
}
\end{ex}

\textbf{Phần 2. Trắc nghiệm đúng sai}
\setcounterpageref{ex}{0}

\begin{ex} %[3D11H59]
Xét tính đúng sai về số có 5 chữ số:
\choiceTF
{\True $10\,000$ là số bé nhất có 5 chữ số.}
{\True $99\,999$ là số lớn nhất có 5 chữ số.}
{Số $100\,000$ có 5 chữ số.}
{\True Số $100\,000$ có 6 chữ số.}
\loigiai{
\begin{itemize}
    \item $10000$ bé nhất có 5 chữ số: Đúng.
    \item $99999$ lớn nhất có 5 chữ số: Đúng.
    \item $100000$ có 6 chữ số (gồm 1 chữ số 1 và 5 chữ số 0). Ý 3 Sai, Ý 4 Đúng.
\end{itemize}
\begin{itemize}
    \item a) \textbf{Đúng}.
    \item b) \textbf{Đúng}.
    \item c) \textbf{Sai}.
    \item d) \textbf{Đúng}.
\end{itemize}
}
\end{ex}

\begin{ex} %[3D11V59]
Cho số $72\,504$. Xét tính đúng sai:
\choiceTF
{\True Chữ số 7 có giá trị là $70\,000$.}
{\True Chữ số 2 có giá trị là $2\,000$.}
{Chữ số hàng chục là chữ số 4.}
{\True Chữ số hàng chục là chữ số 0, hàng đơn vị là 4.}
\loigiai{
\begin{itemize}
    \item Giá trị chữ số 7 là $70000$: Đúng.
    \item Giá trị chữ số 2 là $2000$: Đúng.
    \item Hàng chục là 0 (không phải 4). Ý 3 Sai, Ý 4 Đúng.
\end{itemize}
\begin{itemize}
    \item a) \textbf{Đúng}.
    \item b) \textbf{Đúng}.
    \item c) \textbf{Sai}.
    \item d) \textbf{Đúng}.
\end{itemize}
}
\end{ex}

\textbf{Phần 3. Trả lời ngắn}
\setcounterpageref{ex}{0}

\begin{ex} %[3D11H59]
Viết số liền sau của số $35\,999$.
\shortans{36000}
\loigiai{
\textbf{Đáp số:} $36000$\\
$35999 + 1 = 36000$.
}
\end{ex}

\begin{ex} %[3D11H59]
Viết số gồm $8$ chục nghìn, $4$ nghìn, $2$ trăm và $5$ đơn vị.
\shortans{84205}
\loigiai{
\textbf{Đáp số:} $84205$\\
$84205$.
}
\end{ex}

\begin{ex} %[3D11V59]
Tìm số lớn nhất có 5 chữ số khác nhau lập từ các chữ số $7, 2, 9, 0, 4$.
\shortans{97420}
\loigiai{
\textbf{Đáp số:} $97420$\\
$97420$.
}
\end{ex}

\begin{ex} %[3D11V59]
Tìm số bé nhất có 5 chữ số khác nhau lập từ các chữ số $5, 0, 8, 1, 3$.
\shortans{10358}
\loigiai{
\textbf{Đáp số:} $10358$\\
$10358$.
}
\end{ex}

\begin{ex} %[3D11V59]
Số liền trước của số $80\,000$ là số nào?
\shortans{79999}
\loigiai{
\textbf{Đáp số:} $79999$\\
$80000 - 1 = 79999$.
}
\end{ex}

\begin{ex} %[3D11C59]
Tìm số tròn chục nghìn lớn nhất có 5 chữ số mà tổng các chữ số của nó bằng 9.
\shortans{90000}
\loigiai{
\textbf{Đáp số:} $90000$\\
Số tròn chục nghìn có dạng $\overline{a0000}$. Ta có $a = 9 \Rightarrow 90000$.
}
\end{ex}
